%%%%%%%%%%%%%%%%%%%%%%%%%%%%%%%%%%%%%%%%%%%%%%%%%%%%%%%%%%%%%%
%%%%%%%%%%%%% MA-0561 Algebra Abstracta I %%%%%%%%%%%%%%%%
%%%%%%%%%%%%%%%%%%%%%%%%%%%%%%%%%%%%%%%%%%%%%%%%%%%%%%%%%%%%%%
\newpage
\subsection*{MA-0561 Algebra Abstracta I}
\addcontentsline{toc}{subsection}{MA-0561 Algebra Abstracta I}

\hypertarget{MA0561}{}
\begin{refsection}
\begin{table}[H]
\centering{
\begin{tabular}{|ll|}
\hline
%\multicolumn{2}{|c|}{\textbf{Nivel de virtualidad del curso:} Virtual}\\ \hline
\textbf{Año:} III  & \textbf{Requisitos:} Algebra Lineal II y Teoría de Números \\ 
\textbf{Ciclo:} V  & \textbf{Correquisitos:} Ninguno \\ 
\textbf{Tipo de curso:} Teórico & \textbf{Horas presenciales por semana:}   \\ 
\textbf{Créditos:} 4 & \textbf{Horas de trabajo independiente por semana:}  \\ \hline
\end{tabular}
}
\end{table}



\subsubsection*{Descripción del curso}

Este curso introductorio de Álgebra Abstracta, diseñado para la carrera de Bachillerato y Licenciatura en Matemática, adopta un enfoque distinto al tradicional de iniciar con el estudio de la teoría de grupos. En su lugar, comenzamos con el estudio de la teoría de anillos, aprovechando la familiaridad previa del estudiantado con ejemplos concretos de anillos y módulos, especialmente en el contexto de espacios vectoriales.

El estudiantado explorará los conceptos fundamentales de la teoría de anillos, incluyendo estructuras, homomorfismos y propiedades importantes. Además, se utilizarán ejemplos concretos y previamente conocidos de anillos para facilitar la comprensión y cimentar los conceptos abstractos. 

Hacia la última parte del curso, el enfoque se desplazará hacia la teoría de módulos, permitiendo a los estudiantes entender mejor las estructuras que generalizan los espacios vectoriales.




\subsubsection*{Objetivos}

\textbf{General:} Introducir a la persona estudiante a los conceptos y métodos propios del álgebra abstracta a través del estudio de la teoría de anillos y la teoría de módulos. \\

\textbf{Específicos:}

\begin{enumerate}
\item Aplicar los conceptos y teoremas sobre anillos y módulos en la demostración de resultados y la resolución de problemas teóricos y prácticos.
\item Ejemplificar distintos objetos del álgebra abstracta como tipos de anillos o módulos.
\item Relacionar y contrastar las propiedades algebraicas de las estructuras de anillos y módulos.
\item Indagar en la literatura sobre temas que permitan complementar su formación.
\item Comunicar ideas matemáticas de manera clara y efectiva en forma oral y escrita.
\end{enumerate}

\subsubsection*{Contenidos}

\begin{enumerate}
\item \textbf{Definiciones, ejemplos y propiedades básicas de anillos (1 semana):} 
\begin{enumerate}
    \item Definiciones y ejemplos básicos: $\mathbb{Z}$, $\mathbb{Z}/n\mathbb{Z}$, anillos de polinomios $R[x]$, producto cartesiano de anillos.
    \item Propiedades básicas: consecuencias sencillas de los axiomas, unicidad de neutros e inversos, cancelación aditiva, etc.
    \item Inversos multiplicativos, unidades, divisores de cero.
    \item Dominios enteros y cuerpos.
    \item El Teorema del Binomio para anillos.
    \item Productos cartesianos de anillos.
\end{enumerate}

\item \textbf{Morfismos, ideales y cocientes de anillos (3 semanas):} 
\begin{enumerate}
    \item Subanillos e ideales. Ejemplos y criterios para probar que un subconjunto es un subanillo o un ideal.

    \item Ideales generados por un conjunto, generadores, ideales finitamente generados, ideales principales y anillos de ideales principales. Caracterización de los elementos de un ideal generado por un conjunto.
    \item Sumas y productos de ideales. Asociatividad y distributividad de estas operaciones.
    \item Homomorfismos, epimorfismos, monomorfismos e isomorfismos. El kernel y la imagen. Ejemplos de preservación de propiedades de anillos bajo isomorfismos. (Ver e.g. \cite[p. 72]{Alu21})
    \item Anillos cociente. Ejemplos con $\mathbb{Z}$ y anillos de polinomios.
    \item Los tres Teoremas de Isomorfismos. El Teorema de la Correspondencia entre ideales de un anillo cociente $A/I$ y los ideales de $A$ que contienen a $I$.
    \item Existencia de ideales maximales: todo ideal propio en un anillo está contenido en un ideal maximal.
    \item Ideales primos e ideales maximales. Caracterización de estos en términos de los respectivos anillos cociente.
    \item El Teorema Chino de los Residuos para anillos.
\end{enumerate}

\item \textbf{Teoría de divisibilidad y factorización en anillos conmutativos (3 semanas):} 
\begin{enumerate}
    \item Definición de divisibilidad en un anillo conmutativo. Elementos asociados. Elementos irreducibles y elementos primos.
    \item Relación entre la divisibilidad de elementos y la contención de los correspondientes ideales principales. (ver e.g. \cite[Teorema II.3.2, pág. 136]{Hun80}).
    \item Relación entre los elementos primos e irreducibles con ideales primos y maximales. (ver e.g. \cite[Teorema II.3.4, pág. 136]{Hun80})
    \item Dominios de Factorización Unica. Ejemplos. Dominios de ideales principales y dominios de factorizaición única. Potencias $n$-ésimas y factorización única (ver e.g. \cite[Lema 6.41, pág. 120]{Alu21}).
    \item Anillos y dominios euclideanos. Anillos euclideanos y anillos de ideales principales. 
    \item El máximo común divisor, coprimalidad, y la identidad de Bezout en anillos conmutativos. Existencia del máximo común divisor en dominios de ideales principales y dominios de factorización única (ver e.g. \cite[Teorema II.3.11, pág. 140]{Hun80})
    \item El algoritmo euclideano en dominios euclideanos.
    \item Construcción del cuerpo de fracciones de un dominio entero.
\end{enumerate}

\item \textbf{Anillos de polinomios y factorización (2 semanas):} 
\begin{enumerate}
    \item Repaso de definiciones básicas sobre polinomios.
    \item Propiedades básicas del grado del producto de dos polinomios.
    \item El Algoritmo de la división y el Teorema del Residuo. Anillos de polinomios sobre un cuerpo y dominios euclideanos. Factorización única en anillos de polinomios.
    \item Raíces de polinomios, el Teorema del Factor. Raíces múltiples.
    \item Cocientes de anillos de polinomios. (ver e.g. \cite[Corolario 7.10 y Teorema 7.11, pág. 135]{Alu21})
    \item Irreducibilidad en anillos de polinomios. Criterios de irreducibilidad. 
\end{enumerate}

\item \textbf{Definiciones, ejemplos y propiedades básicas de módulos (2 semanas):} 
\begin{enumerate}
    \item Definición y ejemplos sencillos de $R$-módulos: $k$-espacios vectoriales, anillos, ideales, anillos cociente, $k[x]$-módulos, $R[x]$ como un $R$-módulo.
    \item $\Z$-módulos como grupos abelianos.
    \item Consecuencias sencillas de los axiomas de módulo.
    \item Homomorfismos de $R$-módulos: ejemplos, composición, kernel, imagen e isomorfismo.
    \item Submódulos: ejemplos, submódulo generado por un subconjunto, módulo cíclico.
    \item Sumas directas (finitas) de $R$-módulos: proyecciones e inyecciones canónicas.
    \item Módulos cociente: el epimorfismo de proyección, cokernel.
    \item Los teoremas de isomorfismo.
\end{enumerate}


\item \textbf{Grupos (4 semanas):} 
\begin{enumerate}
    \item Definición y ejemplos.
    \item Homomorfismos y subgrupos.
    \item Potencias de elementos y grupos cíclicos.
    \item Coclases y resultados básicos de conteo. El Teorema de Lagrange.
    \item Normalidad, grupos cociente y homomorfismos. Los teoremas de isomorfismo.

\end{enumerate}
\end{enumerate}



\subsubsection*{Metodología}

\noindent \textbf{Lineamientos metodológicos:} La persona docente a cargo de este curso debe respetar las siguientes pautas:
\begin{enumerate}
    \item Se debe respetar el orden de los temas empezando con anillos y finalizando con módulos.
    \item Fomentar una visión integral de las matemáticas y de cómo se enlazan las diferentes áreas a través del álgebra. Por ejemplo sus relaciones con geometría algebraica, topología algebraica, teoría de números, criptografía, lógica, etc.
\end{enumerate}

\textbf{Sugerencias metodológicas:} Adicionalmente, se tienen las siguientes recomendaciones para la persona docente a cargo de este curso:
\begin{enumerate}
    \item Aprovechar los ejemplos sobre anillos de enteros en cuerpos cuadráticos $\Q(\sqrt{d})$ que fueron estudiados en el curso MA-0496 de Teoría de Números.
    \item Utilizar el tema de módulos para introducir y ejemplificar distintos conceptos de grupos a través de los grupos abelianos vistos como $\Z$-módulos.
    \item Utilizar el recurso tecnológico para reforzar distintos conceptos estudiados en el curso por medio de la visualización, cálculo y experimentación.
    \item Aprovechar momentos adecuados del curso para motivar el estudio por medio de reseñas acerca del desarrollo histórico del álgebra y su importancia en aplicaciones dentro y fuera de la matemática
    \item En la evaluación, se sugiere utilizar estrategias que promuevan el trabajo constante de parte de las personas estudiantes a lo largo del ciclo lectivo. Por ejemplo, esto se puede hacer por medio de tareas o quizzes.
    \item Para fomentar el desarrollo de las habilidades investigativas en el estudiantado, se sugiere la implementación de proyectos de investigación y participación en coloquios o seminarios de investigación.
    \item Promover la comunicación clara y rigurosa de argumentos e ideas matemáticas de manera oral y escrita.
\end{enumerate}


\subsubsection*{Guía para las referencias}

Siempre es bueno conocer distintas referencias pues los diferentes enfoques que se encuentran ayudan a una mejor comprensión de los temas. Aparte del libo de Thomas W. Hungerford \cite{Hun80}, que será nuestra referencia principal en el curso, se recomienda consultar los libros de Gregory T. Lee \cite{Lee18} y de Dan Saracino \cite{Sar08}. El libro de Fernando Q. Gouvêa \cite{Gou12} es una excelente guía que resume una impresionante cantidad de resultados y definiciones básicas de las distintas estructuras básicas del álgebra abstracta, por lo que se recomienda como un libro de repaso una vez finalizado este curso. El libro de Paolo Aluffi \cite{Alu09} aborda el tratamiento del álgebra desde un punto de vista categórico desde un inicio, enfatizando puntos de vista de propiedades universales. Por otro lado, los dos tomos del libro de álgebra moderna avanzada de Joseph Rotman \cite{Rot15} y \cite{Rot17} dan un enfoque algo diferente, empezando el desarrollo de la teoría con un estudio de la teoría de anillos en lugar de la teoría de grupos, esto pues se argumenta que los estudiantes están más familiarizados con ejemplos de anillos que con ejemplos de grupos. El libro de Serge Lang \cite{Lan02} es una gran referencia (aunque quizás no es tan recomendado para aprender la materia por primera vez, por su tratamiento sumamente conciso y seco). Para los estudiantes interesados en aprender Teoría de Números, no se puede hacer algo mejor que recomendar el excelente libro de Kenneth Ireland y Michael Rosen \cite{IR90}. Este libro abarca una amplia gama de temas en la teoría de números, usando un enfoque más algebraico que asume conocimientos de teoría de grupos, anillos y cuerpos a un nivel básico. Para los estudiantes interesados en aplicaciones a la criptografía, el libro de Jeffrey Hoffstein, Jill Pipher y Joseph Silverman \cite{HPS14} es una excelente referencia, que incluso introduce la criptografía con curvas elípticas. También se recomienda el capítulo sobre criptografía del libro de William Stein \cite{Ste09}. Adicionalmente, para aprender sobre el uso de \texttt{SageMath} recomendamos los libros de Razvan Mezei \cite{Mez16} para una introducción rápida. También, se recomienda leer el artículo de William Stein \cite{Ste13} para aprender sobre la filosofía de \texttt{SageMath}, su historia y su relación con otros sistemas CAS. Para temas de historia en general se puede consultar por ejemplo el libro de John Stillwell \cite{Sti10} y más específicamente para historia del álgebra se recomienda consultar el libro de Israel Kleiner \cite{Kle07}.




\nocite{Alu09}
\nocite{Alu21}
\nocite{Lan02}
\nocite{Lan05}
\nocite{Lee18}
\nocite{Hun80}
\nocite{Hun13}
\nocite{AF15}
\nocite{Sti10}
\nocite{Kle07}
\nocite{IR90}
\nocite{Gou12}
\nocite{Kna16a}
\nocite{Kna16b}
\nocite{CE20}


\printbibliography[heading=subbibliography]
\end{refsection}
